% This is samplepaper.tex, a sample chapter demonstrating the
% LLNCS macro package for Springer Computer Science proceedings;
% Version 2.20 of 2017/10/04
%
\documentclass[runningheads]{llncs}
%
\usepackage{graphicx}
\usepackage{comment}
\usepackage{caption}


% Package to generate and customize Algorithm as per ACM style
\usepackage[ruled, linesnumbered, noend]{algorithm2e}
\renewcommand{\algorithmcfname}{Listing}
\SetAlFnt{\small}
\SetAlCapFnt{\small}
\SetAlCapNameFnt{\small}
\SetAlCapHSkip{0pt}
\IncMargin{-\parindent}

%\setlength{\textfloatsep}{0pt}
%\setlength{\intextsep}{0pt}

% Used for displaying a sample figure. If possible, figure files should
% be included in EPS format.
%
% If you use the hyperref package, please uncomment the following line
% to display URLs in blue roman font according to Springer's eBook style:
% \renewcommand\UrlFont{\color{blue}\rmfamily}

\newtheorem{defn}{Definition}

\begin{document}
%
\title{On the Valiant's Algorithm Parallelization\thanks{The research was supported by the Russian Science Foundation, grant No. 18-11-00100.}}
%
%\titlerunning{Abbreviated paper title}
% If the paper title is too long for the running head, you can set
% an abbreviated paper title here
%
\author{Yuliya Susanina\inst{1, 2} \and
Anna Yaveyn\inst{1} \and
Semyon Grigorev\inst{1, 2}\orcidID{0000-0002-7966-0698}}
%
\authorrunning{Yuliya Susanina, Anna Yaveyn, and Semyon Grigorev}
% First names are abbreviated in the running head.
% If there are more than two authors, 'et al.' is used.
%
\institute{Saint Petersburg State University, 7/9 Universitetskaya nab.\\ St. Petersburg, 199034 Russia \\
\and JetBrains Research, Universitetskaya emb., 7-9-11/5A \\ St.Petersburg, Russia \\
\email{jsusanina@gmail.com},
\email{anya.yaveyn@yandex.ru},
\email{semen.grigorev@jetbrains.com}
}
%
\maketitle              % typeset the header of the contribution
%
\begin{abstract}
Recent  research  has  shown  that the theory of formal languages can be used in bioinformatics. While using parsing for processing nucleotide sequences it became necessary to find an easily adaptable to the string-matching problem parsing algorithm, and as such field of application as bioinformatics requires working with large amount of data, it should be highly efficient. The asymptotically fastest parsing algorithm was proposed by Valiant and based on matrix multiplication. The original version of matrix-based algorithm is difficult to apply to the string-matching problem. In this paper we present  a modification  of Valiant's  algorithm dealing with the problem mentioned above. The modification has a succinct proof of correctness and is implemented.

\keywords{Parsing \and Matrix multiplication \and Context-free grammars.}
\end{abstract}
%
%
%
\section{Introduction}

Context-Free Path Querying (CFPQ) is an actively developed area in graph database analysis.
CFPQ is also used for static code analysis~\cite{Reps,10.1145/193173.195287,Zheng}, RDF querying~\cite{10.1007/978-3-319-46523-4_38,MEDEIROS201975}, biological data analysis~\cite{cfpqBio}.

Most of research is focused on developping algorithms for CFPQ evaluation~\cite{hellingsRelational,ward2008distributed,cfpqBio,MEDEIROS201975,Azimov:2018:CPQ:3210259.3210264,Grigorev:2017:CPQ:3166094.3166104}, whereas specification languages for context-free queries are not investigated enough.
Best to our knowledge, only one extension for Sparql supports context-free constarints: cfSPARQL~\cite{10.1007/978-3-319-46523-4_38}.
There is also a proposal for CFPQ as a part of Cypher\footnote{Proposal with path pattern syntax for openCypher: \url{https://github.com/thobe/openCypher/blob/rpq/cip/1.accepted/CIP2017-02-06-Path-Patterns.adoc}.
It is shown that context-free constraints can be expressed with the proposed syntax. Access date: 30.03.2020} language, but there is no implementation for it yet.
We believe that more research should be conducted on the specification languages fo context-free constraints in graph querying.

It is worth noting that graph analysis is often only a part of a more complex system, usually implemented in a general-purpose language.
Since a graph query language is unsuitable to implement a whole system, there should be means of integration of them into general-purpose programming languages.
There are many ways to integrate them ranging from creating graph queries from string values of a general-purpose language to implementing a special embedded domain specific language, and even more sophisticated.

Although simple, the string manipulating approach does not provide a developper with any safety guarantees.
There is no way to ensure that a string generated by an application is a valid query or, in case it is not, to provide any feedback.
This makes string manipulating technique error prone, the code --- unclear and hard to maintain.

Safety of an embedded DSL entirely depends on its implementation.
Some general-purpose languages with powerfull type systems (such as \haskell{}, \ocaml{} or \scala{}) or the ones supporting hygienic macros (such as \scheme{} or \rust{}) facilitate creating safe and reliable DSLs.
Still, they typically lack full support of a development environment: it may be harder to debug queries or issues with composability may arise.

There is a general trend towards imposing more restricting type systems on programming languages.
Among many others are typing annotations for \python{} and \typescript{} code and nullability checks in \kotlin{}.
Typing graphs and query languages improves  readability and simplifies maintainance~\cite{10.1145/2076623.2076653}.

Parser combinators are the answer to the integration of parsing into a general-purpose programming language.
Recursive descend parsers are encoded as functions of the host language, while grammar constructions such as sequencing and choice are implemented as higher-order functions.
This idea was first introduced in~\cite{burge} and further developped in numerous works.
Notable development is monadic parser combinators~\cite{hutton1996monadic}.
In this approach, one can not only parse the input, but simultaneously run semantics calculation if parsing succeeds.
Paper~\cite{izmaylova2016practical} proposed the first monadic parser combinator library which solves the long-standing problem of inability to handle ambiguous and left-recursive grammars.
A library for graph querying was developped~\cite{10.1145/3241653.3241655} based on this work.
The core idea is to use generalized parser combinators as both a way to formulate a query and to execute it.
This approach inherits benefits of combinatory parsing: ease of code reuse, type safety guaranteed by the host language and, since the parser is simply a function, the integrated development support.

Besides integration, it can compute both the single-source and all pairs semantics, as well as execute user actions.
The~single-source semantics is relevant to many real-world applications, including manual data analysis.
It also may be less time-intensive, since on average it needs to expore only a subgraph of the input graph.
Many querying algorithms are only capable to compute all pairs reachability which makes them unsuitable for some applications.

In this paper we make the following contributions.
\begin{itemize}
  \item We demonstrate how to use combinatory-based graph querying on example.
  \item We illustrate such features of the approach as type-safety, flexibility (composability and generics), IDE support and computing user-defined actions.
  \item We evaluate single-source context-free path querying on some real-world RDFs.
  \begin{itemize}
    \item Based on our evaluation, the most common case in RDF context-free querying is when the number of paths in the answer set is big, but they are small.
    \item We demonstrate that the single-source CFPQ can feasibly be used to evaluate such queries.
    \item We conclude that there is a need for a further detailed analysis of both theoretical time and space complexity of single-source CFPQ.
  \end{itemize}
\end{itemize}

\section{Background}

In this section we briefly introduce the key definitions and the necessary parsing algorithm.
Before descibing the Valiant's  algorithm, we would like to mention one of the basic recognition algorithms known as CYK (the Cocke-Younger-Kasami algorithm), by which we show the main Valiant's idea that made such time complexity possible.

\subsection{Preliminaries}

An alphabet $\Sigma$ is a finite nonempty set of symbols. $\Sigma^{*}$ is a set of all finite strings over $\Sigma$.
A grammar is a quadruple $(\Sigma, N, R, S)$, where $\Sigma$ is a finite set of terminals, $N$ is a finite set of nonterminals, $R$ is a finite set of productions of the form $\alpha \rightarrow \gamma$, where $\alpha \in V^{*}NV^{*}$, $\gamma \in V^{*}$, $V = \Sigma \cup N$ and $S \in N$ is a start symbol.

\begin{defn} Grammar $G = (\Sigma, N, R, S)$ is called context-free, if $\space\forall r \in R$ are of the form $A \rightarrow \beta$, where $A \in N, \beta \in V^{+}$.
\end{defn}

\begin{defn} Context-free grammar $G = (\Sigma, N, R, S)$ is said to be in Chomsky normal form if all productions in R are of the form:
\begin{itemize}
  \item $A \rightarrow BC$,
  \item $A \rightarrow a$,
  \item $S \rightarrow \varepsilon$,
\end{itemize}
where $A, B, C \in N, a \in \Sigma, \varepsilon$ is an empty string.
\end{defn}

\begin{defn} $L_{G}(A)$ is language of grammar $G_{A} = (\Sigma, N, R, A)$, which means all the sentences that can be derived in a finite number of rules applications from the start symbol $A$.
\end{defn}


\subsection{Parsing by matrix multiplication}

The main problem of parsing is to verify if the input string belongs to the language of some given grammar $G$. We will describe the Cocke-Younger-Kasami algorithm and the most asymptotically efficient parsing algorithm, which works for all context-free grammars, Valiant's parsing algorithm, based on matrix multiplication. In this paper we use the rewritten version of Valiant's algorithm proposed by Alexander Okhotin.

The CYK algorithm is a basic parsing algorithm.
Its main idea is to construct a parsing table $T$ of size $n \times n$ for an input string $a_{1}a_{2} \dots a_{n}$ and context-free grammar $G = (\Sigma, N, R, S)$ which is in Chomsky normal form, where

$$T_{i, j} =  \{ A | A \in N, a_{i + 1} \dots a_{j} \in L_{G}(A)\} \quad \forall i < j.$$

The elements of $T$ are filled successively beginning with
$ T_{i - 1, i} = \{ A | A \rightarrow a_{i} \in R\}.$
Then, $ T_{i, j} = f(P_{i, j})$ where
$$ P_{i, j} = \bigcup\limits_{k = i + 1}^{j - 1} T_{i,k} \times T_{k, j}  $$
$$ f(P) = \{A | \exists A \rightarrow BC \in R : (B, C) \in P\} $$

The input string $a_{1}a_{2} \dots a_{n}$ belongs to $L_{G}(S)$ if and only if $S \in T_{0, n}$.

% Algorithm1
\begin{algorithm}[t]
\SetAlgoNoLine
\KwIn{Grammar $G = (\Sigma, N, R, S), w = a_{1} \dots a_{n}, n \geq 1, a_{i} \in \Sigma$, where  n + 1 is a power of two}
\underline{main()}{:}{

 \textit{compute(0, n + 1)\;}
 accept if and only if $S \in T_{0, n}$
 \linebreak
 }

\underline{compute(\textit{l, m})}{:}{

 \If {$m - l \geq 4$}{
     \textit{compute(l, $\frac{l+m}{2}$)\;
     compute($\frac{l+m}{2}$, m)}}
 \textit{complete(l, $\frac{l+m}{2}$, $\frac{l+m}{2}$, m)}
 \linebreak
 }

\underline{complete(\textit{l, m}, $l^\prime$, $m^\prime$)}{:}{

 \If {$m - l = 4$ and $m = l^\prime$}{$T_{l, l + 1} = \{A | A \rightarrow a_{l+ 1} \in R\}$\;}
 \ElseIf{$m - l = 1$ and $m < l^\prime$}{ $T_{l, l'} = f(P_{l, l'})$\;}
 \ElseIf{$m - l > 1$}{
    $leftgrounded = (l, \frac{l+m}{2}, \frac{l+m}{2}, m), rightgrounded = (l', \frac{l'+m'}{2}, \frac{l'+m'}{2}, m')$,

    $bottom = (\frac{l+m}{2}, m, l', \frac{l'+m'}{2}), left = (l, \frac{l+m}{2}, l', \frac{l'+m'}{2})$,

    $right = (\frac{l+m}{2}, m, \frac{l'+m'}{2}, m'), top = (l, \frac{l+m}{2}, \frac{l'+m'}{2}, m')$\;
    complete(bottom)\;
    $P_{left} = P_{left} \cup (T_{leftgrounded} \times T_{bottom})$\;
    complete(left)\;
    $P_{right} = P_{right} \cup (T_{bottom} \times T_{rightgrounded})$\;
    complete(right)\;
    $P_{top} = P_{top} \cup (T_{leftgrounded} \times T_{right})$\;
    $P_{top} = P_{top} \cup (T_{left} \times T_{rightgrounded})$\;
    complete(top)
    }
 }
\caption{Parsing by matrix multiplication: Valiant's Version}
\label{algo:valiant}
\end{algorithm}

The time complexity of this algorithm is $O(n^3)$. Valiant proposed to offload the most intensive computations to the Boolean matrix multiplication. As the most time-consuming is computing $\bigcup\limits_{k = i + 1}^{j - 1} T_{i, k} \times T_{k, j}$, Valiant rearranged computation of $T_{i, j}$, in order to use multiplication of submatrices of $T$.

\begin{defn}\label{def:defn4} Let $X \in (2^N)^{m \times l}$ and $Y \in (2^N)^{l \times n}$ be two submatrices of parsing table $T$. Then, $X \times Y = Z$, where $Z \in (2^{N \times N})^{m \times n}$ and $Z_{i, j} = \bigcup\limits_{k = 1}^{l} X_{i, k} \times Y_{k, j}$.
\end{defn}

In listing~\ref{algo:valiant} full pseudo-code of Valiant's algorithm is written in the terms proposed by Okhotin, is presented. All elements of $T$ and $P$ are initialized by empty sets. Then, the elements of these two table are successively filled by two recursive procedures. 

In figure~\ref{fig1} is presented a simple example of Valiant's algorithm. Only the beginning of the work is shown, because later we point out at this version and our approach differences.

\begin{figure}
    \centering
    \includegraphics[width=0.600\textwidth]{pictures/valbeg.pdf}
    \caption{An example of beginning of Valiant's algorithm}
    \label{fig1}
\end{figure}

The procedure $compute(l, m)$ constructs the correct values of $T_{i,j} \forall l \le i < j < m$.

The procedure $complete(l, m, l', m')$ constructs the submatrix $\forall T_{i, j}$ $l \le i < m$, $l' \le j < m'$. This procedure assumes $T_{i, j} \forall l \leq i < j < m,  l' \leq i < j < m'$ are already constructed and the current value of  $P[i, j] =  \{ (B, C) |\exists (m \le k < l'): a_{i + 1} \dots a_{k} \in L(B), a_{k + 1} \dots a_{j} \in L(C)\}$ $\forall l \leq i < m,  l' \leq j < m'$. The submatrix division during the procedure call is shown in figure~\ref{fig2}.

Then Valiant described that product of multiplying of two submatrices of parsing table $T$ can be provided as $|N|^2$ Boolean matrices (for each pair of nonterminals). Denote matrix corresponding to pair $(B, C) \in N \times N$ as $Z^{(B, C)}$, then $Z_{i, j}^{(B, C)} = 1$ if and only if $(B, C) \in Z_{i, j}$. It should also be noted that $Z^{(B, C)} = X^{B} \times Y^{C}$. So, matrix multiplication in definition~\ref{def:defn4} can be replaced by Boolean matrix multiplication, each of which can be computed independently. Following these changes, time complexity of algorithm in listing~\ref{algo:valiant} is $O(|G|BMM(n)log(n))$ for an input string of length $n$, where $BMM(n)$ is the number of operations needed to multiply two Boolean matrices of size $n \times n$.

\begin{figure}
    \centering
    \captionsetup{justification=centering}
    \begin{minipage}{0.40\textwidth}
        \centering
        \includegraphics[width=1\textwidth]{pictures/splitting_with_grounded.pdf}
        \caption{Matrix partition used in \textit{complete(l, m, l', m')} procedure.}
        \label{fig2}
    \end{minipage}\hfill
    \begin{minipage}{0.40\textwidth}
        \centering
        \includegraphics[width=1\textwidth]{pictures/layers.pdf}
        \caption{Matrix partition on V-shaped layers.}
        \label{fig3}
    \end{minipage}
\end{figure}

\section{Modified Valiant's algorithm}

In this section we describe the reorganization of submutrix processing order in the Valiant's algorithm which simplify independent handling of submatrices. As a result, proposed modification simplfy implementation of parallel submatrix processing.

\subsection{New approach}

The main change of this modification is the possibility to divide the parsing table into layers of disjoint submatrices of the same size.
The idea of division we have made from the reorganization of the matrix multiplication order is presented in figure~\ref{fig2}.
Each layer consists of squere matricres which size is power of 2.
The layers are computed successively in the bottom-up order.
Each matrix in the layer can be handled independently, which can help to implement parallel version of layer processing function.

\begin{figure}
\includegraphics[width=250pt]{layers2.eps}
\centering
\caption{Matrix partition on V-shaped layers.}
\label{fig2}
\end{figure}

The modofided viresion of Valiant's algorithm is presented in listing~\ref{algo:modified}.
The procedure \textit{main()} computes the lowest layer $(T_{l, l+1})$, and then divide the table into layers, described earlier, and computes them through the \textit{completeVLayer()} call.
Thus, \textit{main()} computes all elements of parsing table $T$.

For the sake of brevity, we introduce \textit{left(m), right(m), top(m), bottom(m), rightgrounded(m)} and \textit{leftgrounded(m)} functions which returns the nessesary submatrix for matrix $m = (l, m, l', m')$ according to the original Valiant's algorithm (figure~\ref{fig1}).

Also denote some subsidiary functions for matrix layer $M$:
 \begin{itemize}
  \item \textit{bottomsublayer(M)}$ = \{bottom(m) | m \in M \}$,
  \item \textit{lrsublayer(M)}$ = \{\textit{left(m)} | m \in M \} \cup \{\textit{right(m)} | m \in M \}$,
  \item \textit{topsublayer(M)}$ = \{top(m) | m \in M \}$.
\end{itemize}

The procedure \textit{completeVLayer(M)} takes an array of disjoint submatrices $M$ which represents a layer.
For each $m = (l, m, l', m') \in M$ this procedure computes \textit{left(m), right(m), top(m)}.
The procedure assumes that the elements of \textit{bottom(m)} and $T_{i, j}$ for all $i$ and $j$ such that $l \leq i < j < m$ and $  l' \leq i < j < m'$ are already constructed.
Also it is assumed that the current value of
$P_{i, j} =  \{ (B, C) | \exists k, (m \le k < l'), a_{i + 1} \dots a_{k} \in L_G(B), a_{k + 1} \dots a_{j} \in L_G(C)\} $ for all $i$ and $j$ such that $l \leq i < m$ and $l' \leq j < m'$.

The procedure \textit{completeLayer(M)} also takes an array of disjoint submatrices $M$, but unlike the previous one, it computes $T_{i, j}$ for all $(i, j) \in m$.
This procedure requires exactly same assumptions on $T_{i, j}$  and $P_{i, j}$  as in the previous case.

In the other world, \textit{completeVLayer(M)} computes !!!SOME INTUITION HERE!!! \textit{completeLayer(M)} !!!SOME INTUITION HERE!!!

Finally, the procedure \textit{performMultiplication(tasks)}, where \textit{tasks} is an array of a triple of submatrices, perform basic step of algoritm: matrix multiplication. It is worth mentioning that, as distinct from the original algorithm, here $|tasks| \ge 1$ and each task can be computed independently.
So, practical implementation of this procedure can isely involve diferent techniques of parallel array processing, such as OpenMP~\ref{!!!}.

\begin{algorithm}[h!]
\SetAlgoNoLine
\KwIn{$G = (\Sigma, N, R, S), w = a_{1} \dots a_{n}, n \geq 1, n + 1 = 2^p, a_{i} \in \Sigma$ }
\underline{main()}{:}{

 \For {$l \in \{1, \ldots, n \}$}{$T_{l, l + 1} = \{A | A \rightarrow a_{l + 1} \in R\}$}
 \For{$1 \le i < k $}{
 layer = \textit{constructLayer(i)}\;
 \textit{completeVLayer(layer)}
 }
 \BlankLine
 }

\underline{constructLayer(i)}{:}{
 \BlankLine
 $\{B | \exists k \geq 0 : B = (k*2^i, (k+1)*2^i, (k + 1)*2^i, (k+2)*2^i) \}$
 \BlankLine
    }
\underline{completeLayer(M)}{:}{
\BlankLine
\If {$\forall (l, m, l', m') \in M \quad (m - l = 1)$}{\For{$ (l, m, l', m') \in M$}{$T_{l, l'} = f(P_{l, l'})$\;}}
\Else{
\textit{completeLayer(bottomsublayer(M))}\;
\textit{completeVLayer(M)}
}
\BlankLine
}

\underline{comleteVLayer(M)}{:}{
 \BlankLine
 \textit{multiplicationTasks$_1$ = \{left$(m_i)$, leftgrounded$(m_i)$, bottom$(m_i) | m_i \in M \} \cup \linebreak  \{right(subm), bottom(subm), rightgrounded(subm) | subm \in M\}$\;}
 \BlankLine
 multiplicationTask2 = $\{top(subm), leftgrounded(subm), right(subm) | subm \in M\}$\;
 \BlankLine
 multiplicationTask3 = $\{top(subm), left(subm), rightgrounded |subm \in M\}$\;
 \BlankLine
 \textit{performMultiplications(multiplicationTask1)}\;
 \textit{completeLayer(lrsublayer(M))}\;
 \textit{performMultiplications(multiplicationTask2)}\;
 \textit{performMultiplications(multiplicationTask3)}\;
 \textit{completeLayer(topsublayer(M))}

 }
 \BlankLine

 \underline{performMultiplication(tasks)}{:}{\\
 \For{$ (m, m1, m2) \in \textit{tasks}$}{$P_{m} = P_{m} \cup (T_{m1} \times T_{m2})$\;}
 }

\caption{Parsing by matrix multiplication: Modified Version}
\label{algo:modified}
\end{algorithm}


\subsection{Correctness and conplexity}

We provide the proof of correctness and time complexity for the proposed modification in this section.
To do it we should prove correctness of subprocedures \textit{completeLayer} and \textit{completeVLayer}.

\begin{theorem}
Let $M$ be a layer. If for all $(l, m, l', m') \in M$:
\begin{enumerate}
  \item $T_{i, j} = \{ A |  a_{i + 1} \dots a_{j} \in L_G(A)\}$ for all $i$ and $j$ such that $l \leq i < j < m$ and $l' \leq i < j < m'$;
  \item $P_{i, j} =  \{ (B, C) |\exists k, (m \le k < l'): a_{i + 1} \dots a_{k} \in L_G(B), a_{k + 1} \dots a_{j} \in L_G(C)\}$ for all $l \leq i < m$ and $l' \leq j < m'$.
\end{enumerate}

Then the procedure \textit{completeLayer(M)}, returns correctly computed sets of $T_{i, j}$ for all $l \leq i \le m$ and $l' \leq j \le m'$ for all $(l, m, l', m') \in M$.
\end{theorem}


\begin{proof}(Proof by induction on $m - l$.)

Let $(l, m, l', m')$ is a typical element of array $M$.
As far as each element can be handled independantly, we pove statements only for one element of $M$.

\underline{\textbf{Basis:}} $m - l = 1$. There is only one element to compute, and $P_{l, l'} =  \{ (B, C) |  a_{l + 1} \dots a_{l'} \in L(B)L(C)\}$. Further, algorithm computes $f(P[l, l'])$ = \linebreak $\{ A |  a_{l + 1} \dots a_{l'} \in L(A)\}$, so $T[l, l']$ computed correctly.

\underline{\textbf{Inductive step:}} Assume that $(l_1, m_1, l_2, m_2)$ is correctly computed for $m_2 - l_2 = m_1 - l_1 > m - l$.

Let us consider complete \textit{completeLayer(M)}, where $m - l > 1$.

Firstly, consider \textit{completeLayer(bottom(M))}.
Theorem conditions are fulfilled, then this call returns correct sets $T_{i, j}$ for all $(i, j) \in bottom$ (hereinafter is means $\forall (i, j) \in m, \forall m \in \textit{bottom} $).
All submatrices with size $ m_1 - l_1 > m - l $, all previous layers and also \textit{bottom(M)} are correct, so,  \textit{completeVLayer(M)} can be called, and \textit{multiplicationByTask(task1)} adds to each
$ P[i, j] \forall (i, j) \in left(M) $
all pairs
$ \{(B, C) |\exists (\frac{l+m}{2} \le k < l'): a_{i + 1} \dots a_{k} \in L(B), a_{k + 1} \dots a_{j} \in L(C)\} $
 and
$ \forall (i, j) \in right (M)) $
 all pairs
$ \{ (B, C) |\exists (m \le k < \frac{l'+m'}{2}): a_{i + 1} \dots a_{k} \in L(B), a_{k + 1} \dots a_{j} \in L(C)\}$.
Now all the theorem conditions are fulfilled so, it is possible to call $\textit{completeLayer(left } \cup  \ right)$, which returns correct sets $T[i, j] \forall (i, j) \in (left \cup right)$.

Next, \textit{multiplicationByTask(task2)} and \textit{multiplicationByTask(task3)} add to each $P[i, j]$ $\forall (i, j) \in top = \{(l, \frac{l+m}{2}, \frac{l'+m'}{2}, m'))\}$ all pairs $\{(B, C) |\exists (\frac{l+m}{2} \le k < m)$ and $(l' \le k < \frac{l'+m'}{2}) : a_{i + 1} \dots a_{k} \in L(B), a_{k + 1} \dots a_{j} \in L(C)\}$. Now all the theorem conditions are fulfilled so, it is possible to call $\textit{completeLayer(top)}$, which returns correct sets $T[i, j]$ $\forall (i, j) \in top$.

Thus, all $T[i, j]$ $\forall (i, j) \in M$ are computed correctly.
\end{proof}

\begin{theorem}
Let ;M; be a submatrix array. Assume that, $T[i, j] =  \{ A |  a_{i + 1} \dots a_{j} \in L(A)\}$ $\forall l \leq i < j < m,  l' \leq i < j < m'$ and $\forall b1 \leq i < b2$,  $b3 \leq j < b4$, where $(b1, b2, b3, b4) = (\frac{l+m}{2}, m, l', \frac{l'+m'}{2})$, also $P[i, j] =  \{ (B, C) |\exists (m \le k < l'): a_{i + 1} \dots a_{k} \in L(B), a_{k + 1} \dots a_{j} \in L(C)\}$ $\forall l \leq i < m,  l' \leq j < m'$ $\forall (l, m, l', m') \in M$.

Then, the procedure $\textit{completeVLayer(M)}$, returns correctly computed sets of $T[i, j]$ $\forall l \leq i \le m,  l' \leq j \le m'$ $\forall (l, m, l', m') \in M$.
\end{theorem}

\begin{proof}
The proof is similar to the proof of Theorem 1.
\end{proof}

\begin{note}
Function $\textit{costructLayer(i)}$ returns $2^{k - i} - 1$ matrices of size $2^i$.
\end{note}

\begin{lemma}
\
\begin{itemize}
 \item $\forall i \in \{ 1, .., k - 1\}$  $\sum{|layer|}$ for the calls of \textit{completeVLayer(layer)} where $\forall (l, m, l', m') \in layer$ with $m - l = 2^{k - i}$  is exactly $2^{2i - 1} - 2^{i - 1}$;
 \item $\forall i \in \{ 1, .., k - 1\}$ products of submatrices of size $2^{k - i} \times 2^{k - i}$ are calculated exactly $2^{2i - 1} - 2^{i}$
\end{itemize}
\end{lemma}

\begin{proof}
The base case: i = 1. $\textit{completeVLayer(layer)}$ where $\forall (l, m, l', m') \in layer$ with $m - l = 2^{k - 1}$ is called only once in the  $\textit{main()}$ and $|layer| = 1$. So, $2^{2i - 1} - 2^{i - 1} = 2^1 - 2^0 = 1$.

For the induction step, assume that $\forall i \in \{ 1, .., j\}$ $\sum{|layer|}$ for the calls of $\textit{completeVLayer(layer)}$ where $\forall (l, m, l', m') \in layer$ with $m - l = 2^{k - i}$  which is exactly $2^{2i - 1} - 2^{i - 1}$.

Let us consider i = j + 1.

Firstly, it is the call of \textit{completeVLayer(costructLayer(k - i))}, where \textit{costructLayer(i)} returns $2^i - 1$ matrices of size $2^i$. Secondly, \textit{completeVLayer(layer)} is called 3 times for the left, right and top submatrices of size $2^{k - (i - 1)}$. Finally, \textit{completeVLayer(layer)} is called 4 times for the bottom, left, right and top submatrices of size $2^{k - (i - 2)}$, except $2^{i - 2} - 1$ matrices which were already computed.

Then, $\sum{|layer|} = 2^{i} - 1 + 3 \times (2^{2(i - 1) - 1} - 2^{(i - 1) - 1}) + 4 \times (2^{2(i - 2) - 1} - 2^{(i - 2) - 1}) - (2^{i - 2} - 1) = 2^{2i - 1} - 2^{i - 1}$.

To calculate the number of products of submatrices of size $2^{k - i} \times 2^{k - i}$, we consider the calls of \textit{completeVLayer(layer)} where $\forall (l, m, l', m') \in layer$ with $m - l = 2^{k - (i - 1)}$, which is $2^{2(i - 1) - 1} - 2^{(i - 1) - 1}$. During these calls \textit{performMultiplications} run 3 times, $|multiplicationTask1| = 2 \times 2^{2(i - 1) - 1} - 2^{(i - 1) - 1}$ and \linebreak $|multiplicationTask2|$ = $|multiplicationTask3| = 2^{2(i - 1) - 1} - 2^{(i - 1) - 1}$. So, the number of products of submatrices of size $2^{k - i} \times 2^{k - i}$ is $4 \times (2^{2(i - 1) - 1} - 2^{(i - 1) - 1}) = 2^{2i - 1} - 2^{i}$.
\end{proof}

\begin{theorem}
The time complexity of the Algorithm 1 is $O(|G|BMM(n)\log{n})$ for an input string of length n, where G is a context-free grammar in Chomsky normal form, BMM(n) is the number of operations needed to multiply two Boolean matrices of size $n \times n$.
\end{theorem}

\begin{proof}
The proof is almost identical with that of the theorem given by Okhotin~\cite{okhotin}, because, as shown in the last lemma, the Algorithm 1 has the same number of products of submatrices.
\end{proof}

To summarize, the correctness of the modification was proved and it was shown that the time complexity remained the same as in Valiant's version.

\subsection{Algorithm for substrings}

Next we show how our modification can be applied to the string-matching problem.

So if we want to find all substrings of size $s$ which can be derived from a start symbol for an input string of size $n = 2^k$, we need to compute layers with submatrices of size not greater than $2^{l'}$, where $2^{l' - 2} < s \le 2^{l' - 1}$.

$l' = k - (m - 2)$, $(m - 2) = k - l'$

$ C \sum\limits_{i=m}^k 2^{2i - 1} \cdot 2^{\omega(k - i)} \cdot f(2^{k - i}) = C \cdot 2^{\omega l'}\sum\limits_{i=2}^{l'} 2^{(2 - \omega)i} \cdot 2^{2(k - l') - 1} \cdot f(2^{l' - i}) \le C \cdot 2^{\omega l'} f(2^{l'}) \cdot 2^{2(k - l') - 1} \sum\limits_{i=2}^{l'} 2^{(2 - \omega)i} = BMM(2^{l'}) \cdot 2^{2(k - l') - 1} \sum\limits_{i=2}^{l'} 2^{(2 - \omega)i}$

Thus, time complexity for searching all substrings is  $O(|G|BMM(2^{l'})(l' - 1))$, while time complexity for the full input string is $O(|G|BMM(2^k)(k - 1))$. In contract to the modification, Valiant's algorithm completely calculate at least 2 triangle submatrices of size $\frac{n}{2}$, which mean minimum asymptotic complexity  $O(|G|BMM(2^{k - 1})(k - 2))$. Make a conclusion that the modification is asymptotically faster for substrings of size $s \ll n$  than the original algorithm.

\section{Evaluation}

The goal of this evaluation is to assess the performance scaling of Spla on Vortex.
Due to limitations in atomic operation support within the RTL implementation, all experiments were performed using the SimX functional simulator.

\subsection{Environment}

Initial testing revealed issues with floating-point operations, which produced incorrect results for some hardware configurations.
Consequently, we limited subsequent experiments to Breadth-First Search (BFS) and Triangle Counting (TC), excluding Single-Source Shortest Path (SSSP) and PageRank.
To keep simulation times manageable, we used a single graph from the SuiteSparse matrix collection\footnote{A diverse collection of sparse matrices from various domains: \url{http://sparse.tamu.edu/}}: soc-Epinions1, with 75~888 vertices and 508~837 edges.


We conducted two series of experiments.
The first varies the number of warps and threads per warp while keeping the number of clusters and cores fixed (at 2 and 4, respectively), with the goal of selecting the best core configuration while preserving multi-core execution to account for cache effects.
The second series, using the best configuration identified in the first step, varies the number of clusters and cores per cluster to assess scaling at the core and cluster levels.
Cache sizes were set to their default values: 16 KB for $L_1$, 1 MB for $L_2$, and 2 MB for $L_3$.

We use the number of cycles reported by SimX as a performance metric.
For multi-core configurations, we report the maximum cycle count across all cores.
During the experiments, we encountered unexpected behavior in SimX that led to out-of-memory exceptions. 
Therefore, some data points are missing from the graphs below.

\subsection{Results}

In figures~\ref{fig:tc_threads_warps} and~\ref{fig:bfs_threads_warps}

\begin{figure}
    \begin{center}
        \includegraphics[width=0.49\textwidth]{pictures/TC_threads_warps.pdf}
    \end{center}
    \caption{Scaling analysis of triangle counting for varying numbers of warps and threads per warp}
    \label{fig:tc_threads_warps}
\end{figure}

\begin{figure}
    \begin{center}
        \includegraphics[width=0.49\textwidth]{pictures/BFS_threads_warps.pdf}
    \end{center}
    \caption{Scaling analysis of BFS for varying numbers of warps and threads per warp}
    \label{fig:bfs_threads_warps}
\end{figure}

Best configuration for BFS is 2 warps, 8 threads per warp (16 threads total). 
Best configuration for TC is 4 warps, 16 threads per warp (64 threads total).


\begin{figure}
    \begin{center}
        \includegraphics[width=0.49\textwidth]{pictures/BFS_cores_clusters.pdf}
    \end{center}
    \caption{Scaling analysis of BFS for varying numbers of clusters and cores per cluster}
    \label{fig:bfs_cores_clusters}
\end{figure}


Edges per core on cycle. Compare with Spla on other GPUs.

\subsection{Scaling limitations analysis}

%sum(scoreboard stalls * lsu_percent) / sum(instr) * 100
To analyze the reasons for limited scaling as the number of threads increases, we measured the average utilization of the ALU and LSU, in terms of stall cycles, for the best BFS configuration.
The results are presented in Fig.~\ref{fig:bfs_alu_stalls} and Fig.~\ref{fig:bfs_lsu_stalls}, respectively.
The data indicate that the LSU is the performance bottleneck within the core.

The same bottleneck was observed in the scaling analysis across clusters and cores.
Whether increasing cache sizes can alleviate this problem remains a question for future research.
We anticipate that careful cache size tuning may help identify a more efficient configuration.

\begin{figure}
    \begin{center}
        \includegraphics[width=0.49\textwidth]{pictures/BFS_alu.pdf}
    \end{center}
    \caption{ALU stalls on BFS for the best configuration}
    \label{fig:bfs_alu_stalls}
\end{figure}

\begin{figure}
    \begin{center}
        \includegraphics[width=0.49\textwidth]{pictures/BFS_lsu.pdf}
    \end{center}
    \caption{LSU stalls on BFS for the best configuration}
    \label{fig:bfs_lsu_stalls}
\end{figure}
\section{Conclusion and Future Work}

Platform presented.

Education. Metaprogramming, translators development, GPGPU programming, etc.

Graph parsing.

Geterogenious porgramming generalization. Hopac is better then MBP~\footnote{\url{https://vasily-kirichenko.github.io/fsharpblog/actors}}.

Research: Automatic memory management.

Data to code translation (automata can be translated into code instead of data structures in memory)

Other technical improvements: IDE support, type provider improvements, new OpenCL standard support, runtime extension, etc.


%
% ---- Bibliography ----
%
% BibTeX users should specify bibliography style 'splncs04'.
% References will then be sorted and formatted in the correct style.
%
\bibliographystyle{splncs04}
\bibliography{Parsing_by_Matrix_Multiplication_for_String_Matching}

\end{document}
